\documentclass[11pt,a4paper]{article}
\usepackage[utf8]{inputenc}
\usepackage[T1]{fontenc}
\usepackage{tgpagella}
\usepackage[spanish]{babel}
\usepackage{geometry}
\usepackage{booktabs}
\usepackage{xcolor}
\usepackage{titlesec}
\usepackage{fancyhdr}
\usepackage[parfill]{parskip}
\usepackage[expansion=false]{microtype}
\usepackage{hyperref}
\usepackage{setspace}
\usepackage{needspace}
\usepackage{etoolbox}
\usepackage{graphicx}

\widowpenalty=10000
\clubpenalty=10000
\displaywidowpenalty=10000
\raggedbottom
\sloppy
\hbadness=10000
\hfuzz=2pt

\geometry{top=3.0cm,bottom=3.0cm,left=3.5cm,right=3.5cm}
\setlength{\parskip}{0.65em}
\setlength{\parindent}{0em}
\setlength{\headheight}{14pt}

\definecolor{azulai}{RGB}{30,80,140}
\definecolor{doradoai}{RGB}{140,90,0}
\definecolor{grisai}{RGB}{70,70,70}

\titleformat{\section}{\Large\bfseries\color{azulai}}{}{0em}{}[{\color{azulai}\titlerule[0.5pt]}]
\titlespacing*{\section}{0pt}{1.8em}{0.8em}

\titleformat{\subsection}{\large\bfseries\color{doradoai}}{}{0em}{}
\titlespacing*{\subsection}{0pt}{1.4em}{0.5em}

\titleformat{\subsubsection}{\normalsize\bfseries\color{grisai}}{}{0em}{}
\titlespacing*{\subsubsection}{0pt}{1.0em}{0.3em}

% Evita títulos huérfanos al final de página
\pretocmd{\section}{\Needspace{12\baselineskip}}{}{}
\pretocmd{\subsection}{\Needspace{9\baselineskip}}{}{}
\pretocmd{\subsubsection}{\Needspace{6\baselineskip}}{}{}

\pagestyle{fancy}
\fancyhf{}
\rhead{\textcolor{grisai}{\small Nieves Arrarás — Arquitectura Interna}}
\lhead{\textcolor{grisai}{\small Luna · Casa 4 · Casa 6}}
\cfoot{\textcolor{grisai}{\small\thepage}}
\renewcommand{\headrulewidth}{0.3pt}

\hypersetup{colorlinks=true,linkcolor=azulai,urlcolor=azulai}
\setstretch{1.45}
\tolerance=2500
\emergencystretch=8em

\begin{document}

% ── Portada ──────────────────────────────────────────────────────────────────
\begin{titlepage}
  \centering
  \vspace*{1.5cm}
  {\Huge\bfseries\color{azulai} Luna · Casa 4 · Casa 6}\\[0.5cm]
  {\large\color{grisai} Arquitectura Interna}\\[0.3cm]
  {\small\itshape\color{grisai} Un método para sostener cuerpo, energía y vida con coherencia}\\[2cm]
  {\huge\color{doradoai} Nieves Arrarás}\\[1.5cm]
  {\Large 01/12/1948 \quad 01:20}\\[0.3cm]
  {\Large Pamplona, España}\\[0.3cm]
  {\normalsize Lat: 42.8157° \quad Lon: -1.6522° \quad UTC+01:00}\\[0.3cm]
  {\normalsize Ascendente: Virgo 16°53' \quad
    MC: Géminis 14°27'}\\[2cm]

  \begin{tabular}{lll}
    \textbf{Luna:} & Sagitario & Casa 3 \\
    \textbf{Casa 4:} & Sagitario en cúspide & Regente: Júpiter \\
    \textbf{Casa 6:} & Acuario en cúspide & Regente: Urano \\
  \end{tabular}\\[2cm]

  \vfill
  {\small Generado el 13/05/2026}
\end{titlepage}

\tableofcontents
\clearpage

% ── Datos de referencia ───────────────────────────────────────────────────────
\section{Datos de referencia}

\begin{center}
\begin{tabular}{llll}
  \toprule
  \textbf{Eje} & \textbf{Signo} & \textbf{Posición} & \textbf{Datos} \\
  \midrule
  Luna & Sagitario & Casa 3 & 11°39' \\
  Casa 4 (cúspide) & Sagitario & --- & Regente: Júpiter \\
  Casa 6 (cúspide) & Acuario & --- & Regente: Urano \\
  Ascendente & Virgo & --- & 16°53' \\
  Medio Cielo & Géminis & --- & 14°27' \\
  \bottomrule
\end{tabular}
\end{center}

\vspace{0.4cm}
\textbf{Planetas en Casa 4:} Marte, Júpiter\\
\textbf{Planetas en Casa 6:} Lilith

\vspace{0.5cm}
\textbf{Aspectos de la Luna:}

\begin{center}
\begin{tabular}{lllll}
  \toprule
  & \textbf{Asp.} & \textbf{Planeta} & \textbf{Tipo} & \textbf{Orbe} \\
  \midrule
  Luna & ✶ & Neptuno & Sextil & 2.9° \\
  Luna & = & Sol & Conjunción & 2.9° \\
  Luna & □ & Lilith & Cuadratura & 3.2° \\
  Luna & △ & Plutón & Trígono & 4.9° \\
  Luna & □ & Saturno & Cuadratura & 5.7° \\
  \bottomrule
\end{tabular}
\end{center}

\vspace{0.8cm}
\begin{center}
\includegraphics[width=0.80\textwidth]{C:/Users/Izas/claude/astrología/pruebas/Nieves_Arrarás_Luna_C4_C6_rueda.png}
\end{center}

\clearpage

% ── Interpretación ────────────────────────────────────────────────────────────
\section{Interpretación — Arquitectura Interna}

\Needspace{8\baselineskip}
\begin{center}
{\small\itshape
No se trata de explicarte. Se trata de mostrar cómo funciona tu regulación interna,\\
dónde puede haber desregulación y qué necesitas para sostenerte.
}
\end{center}

\vspace{0.5cm}

% ── 1. Luna ───────────────────────────────────────────────────────────────────
\subsection{1. Tu Luna — Función emocional}

\subsubsection*{Luna en Sagitario — Casa 3}

Necesitas sentir que lo que estás viviendo tiene dirección.

Cuando una emoción encuentra sentido o movimiento, puedes atravesarla mucho mejor. El problema aparece cuando sientes algo que no sabes hacia dónde va o que parece no conducir a ningún lugar.

Entonces suele aparecer inquietud, necesidad de escapar, cambiar de aire o salir de la situación de alguna manera. Tu cuerpo necesita espacio, movimiento y sensación de horizonte.

Cuando pasas demasiado tiempo sintiendo que no existe libertad o amplitud suficiente, la activación empieza a acumularse.

La regulación mejora cuando existe dirección, expansión y posibilidad de avanzar. En cambio, la sensación de límite, encierro o falta de sentido tiende a aumentar rápidamente la inquietud interna.

Necesitas poner palabras a lo que sientes para poder ordenarlo internamente.

Cuando algo no puede hablarse, escribirse o expresarse, se queda circulando dentro de tu cabeza.

Muchas veces la emoción se transforma entonces en pensamiento repetitivo, ruido mental o necesidad constante de entender lo que está ocurriendo.

El entorno cotidiano tiene un impacto importante sobre tu estado emocional: las conversaciones, los mensajes, los intercambios y los vínculos cercanos te afectan más de lo que suele parecer.

Cuando acumulas demasiadas cosas sin expresarlas, la mente rara vez consigue descansar del todo.

La regulación mejora cuando puedes hablar, escribir o sentir que alguien realmente te escucha. En cambio, el silencio sostenido, guardarte lo que te pasa o no poder expresarlo tiende a aumentar la saturación interna.

Existe una vía suave entre tu mundo emocional y tu sensibilidad más sutil.

La imaginación, el descanso, la música, el silencio o la contemplación pueden ayudarte a regularte. Cuando cuidas tus límites, la sensibilidad deja de sentirse como sobrecarga.

La regulación mejora cuando hay apertura, pero también contorno.

Lo que sientes y la dirección vital suelen moverse muy cerca.

Cuando emocionalmente hay estabilidad, aparece más energía, claridad y sensación de propósito. Y cuando algo se rompe por dentro, la fuerza vital también lo nota enseguida.

No separas fácilmente quién eres de cómo te sientes en ese momento. Eso puede darte mucha coherencia interna, pero también hacer que los períodos difíciles afecten con más profundidad.

Puede haber tensión entre tus necesidades emocionales y una parte instintiva que se resiste a la adaptación.

Cuando aparece sensación de invasión, control o condicionamiento emocional, puede surgir una reacción fuerte. Esa señal tiene información importante, aunque no siempre necesite convertirse en corte.

El trabajo es escucharla sin dejar que destruya automáticamente el vínculo o la situación.

Tienes una capacidad profunda de transformación emocional.

Cuando algo te afecta, puedes atravesarlo con una fuerza interna muy grande. No siempre es fácil, pero existe capacidad real para regenerarte desde dentro.

La intensidad, cuando encuentra cauce, puede convertirse en poder de transformación.

Puede costarte permitir ciertas emociones sin intentar controlarlas enseguida.

La vulnerabilidad puede sentirse peligrosa, excesiva o difícil de mostrar. Muchas veces una parte de ti intenta contener lo que sientes antes incluso de haberlo reconocido completamente.

No necesitas destruir tu estructura. Necesitas permitir que la emoción exista dentro de ella.

% ── 2. Casa 4 ─────────────────────────────────────────────────────────────────
\subsection{2. Tu Casa 4 — Raíz interna}

\subsubsection*{Sagitario en la cúspide — Regente: Júpiter}

Necesitas sentir que tu vida tiene dirección para experimentar estabilidad interna.

Cuando existe horizonte, expansión o sensación de crecimiento, tu mundo interno se sostiene mucho mejor.

La dificultad aparece cuando la vida empieza a sentirse demasiado estrecha, repetitiva o sin sentido. Entonces puede surgir inquietud, necesidad de escapar o sensación de encierro interno.

Muchas veces intentas regularte tomando distancia, moviéndote o buscando aire y perspectiva.

La regulación mejora cuando existe libertad, movimiento y sensación de propósito. En cambio, sentir límites constantes, falta de salida o vivir algo que no tiene sentido para ti tiende a aumentar rápidamente la activación.

El regente de la Casa 4 es Júpiter, en Capricornio y en Casa 4. La base interna tiende a organizarse desde las raíces y el hogar interno, con la cualidad concreta y sostenida de Capricornio.

Cuando algo amenaza tu espacio emocional o tu sensación de seguridad, la reacción suele aparecer muy rápido.

El cuerpo entra fácilmente en estado de defensa cuando percibe invasión, tensión o conflicto dentro del entorno íntimo.

Muchas veces la respuesta llega antes incluso de haber comprobado completamente si el peligro es real.

La sensación de tener un espacio propio y protegido es especialmente importante para ti.

Cuando no existe esa sensación de resguardo, el nivel de activación interna sube rápidamente.

Tu mundo interior vive las experiencias a gran escala.

Cuando algo toca tu base emocional, la experiencia puede expandirse muchísimo por dentro.

Existe una necesidad muy profunda de encontrar un lugar que realmente se sienta hogar.

Y muchas veces ningún espacio externo parece cubrir completamente esa necesidad.

Cuando aparece sensación de inseguridad emocional, la tendencia suele ser expandirte: buscar más sentido, más espacio o nuevas posibilidades.

Necesitas amplitud también dentro de tu vida íntima.

% ── 3. Casa 6 ─────────────────────────────────────────────────────────────────
\subsection{3. Tu Casa 6 — Regulación cotidiana}

\subsubsection*{Acuario en la cúspide — Regente: Urano}

Necesitas hacer las cosas a tu manera para poder sostenerlas.

Las rutinas demasiado rígidas o impuestas desde fuera suelen agotarte rápidamente.

Tu sistema nervioso es especialmente sensible al exceso de estímulos, ruido o demandas constantes.

Necesitas períodos reales de desconexión para volver a sentir claridad interna.

Cuando pasas demasiado tiempo adaptándote a ritmos externos que no se sienten propios, el cuerpo empieza a saturarse.

Entender por qué haces algo tiene muchísimo impacto sobre tu capacidad para sostenerlo en el tiempo.

La regulación mejora cuando existe autonomía, espacio mental y libertad de ritmo. En cambio, la sobreestimulación, la imposición externa o no disponer de tiempo para desconectar tienden a generar mucha saturación nerviosa.

El regente de la Casa 6 es Urano, en Géminis y en Casa 10. La regulación cotidiana tiende a organizarse desde la vocación y el espacio público, con la cualidad mental y relacional de Géminis.

Tu cuerpo no siempre responde bien a normas impuestas desde fuera.

Existe una parte muy instintiva en tu manera de relacionarte con la salud, el descanso y el cuidado cotidiano.

Muchas veces necesitas descubrir personalmente qué es lo que realmente te funciona.

Cuando intentas forzarte continuamente a seguir modelos que no encajan contigo, el cuerpo acaba reaccionando.

Escuchar lo que necesitas de verdad suele ser mucho más importante que intentar hacer lo correcto según otras personas.

% ── 4. Integración ────────────────────────────────────────────────────────────
\subsection{4. Integración — Cómo funcionan las tres capas juntas}

Tu Luna en Sagitario muestra cómo reaccionas emocionalmente cuando algo te afecta.

Es el primer movimiento: lo que aparece antes de que puedas ordenarlo, explicarlo o decidir qué hacer con ello.

Después, esa emoción busca un lugar donde apoyarse. Ahí entra tu Casa 4, con Sagitario en la cúspide: la parte de ti que necesita sentir base, refugio y seguridad interna.

Si esa base está disponible, lo que sientes puede moverse con más facilidad.
Si no lo está, la emoción se queda más tiempo en estado de alerta.

Y lo que no encuentra apoyo interno termina llegando al cuerpo. Ahí entra tu Casa 6, con Acuario en la cúspide: tus hábitos, tu energía diaria, tu descanso y la forma en que el cuerpo sostiene lo que vives.

Por eso, cuando algo se desordena emocionalmente, no se queda solo en lo emocional. Puede acabar apareciendo como cansancio, pérdida de ritmo, dificultad para descansar o sensación de no poder sostener el día con la misma energía.

Tu Luna y tu Casa 4 comparten el elemento Fuego.

Eso significa que lo que sientes y lo que necesitas para sentir base interna tienen un lenguaje parecido.

Cuando emocionalmente algo se activa, una parte profunda de ti reconoce bastante rápido qué necesita para sentirse más segura.

La Casa 6, en cambio, está en Aire. Tu cuerpo y tus hábitos funcionan con otra lógica.

Puede que entiendas emocionalmente lo que necesitas, pero aun así te cueste llevarlo al ritmo cotidiano, al descanso, al cuerpo o a la rutina.

Además, Marte y Júpiter en la Casa 4 influye directamente en tu sensación de base interna.

Cuando buscas seguridad, refugio o estabilidad emocional profunda, esa energía ya está presente ahí.

Por eso tu forma de sentir hogar, intimidad y raíz no depende solo del signo de la Casa 4. También está marcada por lo que esos planetas activan dentro de ti.

Además, Lilith en la Casa 6 influye en cómo tu cuerpo sostiene la vida cotidiana.

No se trata solo de hábitos o rutina. Se trata de cómo organizas tu energía, cómo respondes al cansancio y qué ocurre cuando el cuerpo empieza a pedir ajuste.

Esos planetas muestran qué tipo de fuerza, sensibilidad o tensión entra en tu regulación diaria.

Los aspectos de tensión con Lilith y Saturno hacen que el mundo emocional tenga menos margen en ciertos momentos.

Cuando algo te activa, puede haber menos tiempo entre sentir y reaccionar. La emoción puede llegar más intensa, más mezclada o más difícil de ordenar.

Esto no significa que sí o sí vayas a desregularte. Significa que necesitas reconocer antes las señales iniciales, porque cuando la activación ya ha subido mucho, cuesta más volver al centro.

Al mismo tiempo, los aspectos de apoyo con Neptuno, el Sol y Plutón ofrecen recursos reales.

Cuando consigues encontrar un primer punto de estabilidad, esos aspectos ayudan a sostenerlo.

No eliminan la tensión, pero pueden facilitar que vuelvas a organizarte sin quedarte completamente dentro de lo que se activó.

Lo más difícil no suele ser que una sola parte se altere.

Lo más difícil aparece cuando las tres capas se mueven al mismo tiempo:
la Luna en Sagitario, con su forma concreta de sentir;
la Casa 4 en Sagitario, con su necesidad de base y seguridad interna;
y la Casa 6 en Acuario, con el cuerpo intentando sostenerlo todo en la vida diaria.

Cuando esto ocurre, puedes notar que no sabes por dónde empezar: emocionalmente hay activación, internamente falta suelo y el cuerpo empieza a mostrar cansancio, tensión o pérdida de ritmo.

El trabajo no consiste en resolverlo todo a la vez.

Consiste en encontrar la primera capa disponible: una emoción que pueda expresarse, un entorno que pueda darte más seguridad, o una acción corporal sencilla que ayude al sistema a volver poco a poco.

% ── 5. Sostén y desregulación ─────────────────────────────────────────────────
\subsection{5. Sostén y desregulación}

Tu Luna en Sagitario se desregula cuando lo que sientes no encuentra salida.

Puede ser porque no puedes expresarlo, porque no hay espacio para moverte, porque el entorno no lo recibe o porque intentas sostenerlo demasiado tiempo dentro.

No siempre hace falta que ocurra algo muy intenso.

A veces basta con acumular durante demasiado tiempo una emoción que no ha tenido lugar, palabra, movimiento o descanso.

Los aspectos de tensión con Lilith y Saturno hacen que ese margen sea más estrecho.

Cuando algo te activa, puede haber menos tiempo entre sentir y reaccionar.

La emoción puede llegar más fuerte, más mezclada o más difícil de ordenar. Por eso es importante reconocer las primeras señales antes de llegar al punto de desborde.

Tu Casa 4 con Sagitario pierde estabilidad cuando dejas de sentir base interna.

No siempre tiene que haber un conflicto claro.

A veces basta con que cambie el clima emocional del entorno, con que una relación cercana se tense, con que no puedas descansar de verdad o con que el espacio íntimo deje de sentirse seguro.

Cuando eso ocurre, no solo cambia el ánimo. Cambia la sensación de tener suelo por dentro.

Tu Casa 6 con Acuario muestra la desregulación en el cuerpo y en la vida cotidiana.

Puede aparecer como cansancio, pérdida de ritmo, tensión física, sueño menos reparador o dificultad para sostener hábitos que normalmente sí puedes sostener.

Muchas veces el cuerpo avisa antes de que puedas explicar qué está pasando.

Cuando el cuerpo empieza a perder continuidad, conviene mirar también qué emoción no ha encontrado salida y qué parte de tu base interna ha dejado de sentirse segura.

Los momentos más difíciles aparecen cuando las tres capas se mueven a la vez.

La Luna en Sagitario se activa.
La Casa 4 en Sagitario deja de sentirse como base.
Y la Casa 6 en Acuario empieza a mostrarlo en el cuerpo, el descanso o la rutina.

Ahí puedes sentir que no sabes por dónde empezar.

No se trata de resolverlo todo a la vez.
Se trata de encontrar la primera capa disponible: expresar algo, descansar, ordenar una pequeña rutina, moverte, pedir espacio o recuperar una mínima sensación de seguridad.

Los aspectos de apoyo con Neptuno y el Sol funcionan como recursos cuando consigues recuperar un primer punto de estabilidad.

No eliminan la desregulación, pero pueden ayudarte a reorganizarte antes.

Cuando aparece una mínima base, esos aspectos facilitan que no te quedes dentro de lo que se activó.

% ── 6. Orientación final ──────────────────────────────────────────────────────
\subsection{6. Orientación final}

\subsubsection*{Qué observar}

Empieza observando qué situaciones activan tu Luna en Sagitario.

No para corregirte.
No para controlar lo que sientes.

Solo para reconocer antes cómo empieza la activación.

¿Qué ocurre justo antes de sentirte que la situación te sobrepasa?
¿Qué cambia en el entorno, en los vínculos o en el ritmo diario?

Muchas veces la desregulación no empieza en el momento del desborde. Empieza bastante antes: cuando tu Casa 4 en Sagitario deja de sentirse segura o cuando tu Casa 6 en Acuario pierde continuidad y el cuerpo empieza a acumular tensión.

\subsubsection*{Qué puede desregularte}

En tu caso, la desregulación suele seguir una secuencia.

Primero, tu Luna en Sagitario busca algo que necesita para estabilizarse y no lo encuentra.

Después, tu Casa 4 deja de funcionar como base interna. Puede ser por tensión emocional, por cambios en el entorno, por exceso de carga o simplemente porque llevas demasiado tiempo sosteniendo más de lo que puedes absorber.

Y finalmente, la Casa 6 en Acuario empieza a mostrarlo en el cuerpo:
menos energía,
menos descanso,
menos capacidad de sostener hábitos o continuidad.

Cuando las tres capas coinciden, puedes sentir que todo cuesta más de lo habitual aunque desde fuera no parezca estar ocurriendo nada tan grave.

\subsubsection*{Qué te estructura y sostiene}

Lo que te ayuda a regularte no tiene que ser enorme.

De hecho, muchas veces lo más importante son cosas pequeñas que se sostienen en el tiempo.

Una mínima rutina corporal.
Un espacio de descanso real.
Un entorno que se sienta más seguro.
Una forma de expresar lo que ocurre antes de acumularlo demasiado.

Cuando tu Casa 4 siente suelo, tu Luna necesita menos esfuerzo para estabilizarse.

Y cuando tu Casa 6 mantiene cierta continuidad —aunque sea pequeña— el cuerpo puede absorber mejor lo que emocionalmente estás atravesando.

No se trata de no sentir.
Se trata de que lo que sientes no tenga que terminar descargándose siempre en agotamiento, tensión o pérdida de equilibrio.

\subsubsection*{Práctica}

Tu regulación necesita expresión mental y espacio para ordenar lo que ocurre.

Reserva unos minutos al día para escribir, hablar o simplemente nombrar tu estado interno.

No hace falta entenderlo todo.
Hace falta que no se quede circulando continuamente dentro.

Cuando algo encuentra palabras, el sistema deja de sostener tanta carga en silencio.

\Needspace{8\baselineskip}
\vspace{0.6cm}

\begin{center}
{\small\itshape\color{grisai}
La astrología se usa aquí como lenguaje simbólico de observación, no como definición de la persona.\\
Lo que tienes delante es un mapa de posibilidades, no un destino fijo.
}
\end{center}

\end{document}
